\section{Dokumentasi Implementasi dan Kode Sumber}
\label{sec:dokumentasi_kode}

Seluruh rangkaian eksperimen dalam penelitian ini diimplementasikan menggunakan bahasa pemrograman Python 3.13 dalam lingkungan pengembangan Jupyter Notebook. Pendekatan ini dipilih untuk mendukung reprodusibilitas (\textit{reproducibility}) dan visualisasi data interaktif. Struktur kode dibagi menjadi beberapa modul notebook terpisah untuk memfasilitasi pengembangan, pengujian, dan analisis perbandingan antar strategi.

Berikut adalah rincian fungsionalitas dari setiap artefak komputasi yang dihasilkan selama proses penelitian:

\subsection{Modul Utama dan Perbandingan}
Bagian ini berisi kode inti yang digunakan untuk menghasilkan hasil akhir tesis, mencakup implementasi strategi \textit{Temporal Graph-Enhanced} dan perbandingannya dengan \textit{Baseline}.

\begin{enumerate}
    \item \textbf{\texttt{comparison\_baseline\_vs\_temporal\_graph\_fixed.ipynb}} \\
    Notebook ini adalah \textbf{artefak utama} penelitian. Notebook ini mengintegrasikan seluruh komponen sistem:
    \begin{itemize}
        \item Konstruksi \textit{Temporal Graph} dari data historis harga cryptocurrency.
        \item Ekstraksi fitur jaringan (\textit{Degree Centrality, Betweenness Centrality, Density}).
        \item Pelatihan model \textit{Machine Learning} (Random Forest) untuk deteksi rezim pasar.
        \item Simulasi \textit{backtesting} strategi Baseline vs Temporal Graph.
        \item Perhitungan metrik kinerja (Sharpe Ratio, Max Drawdown, Return) dan uji signifikansi.
    \end{itemize}
    
    \item \textbf{\texttt{temporal\_graph\_markowitz.ipynb}} \\
    Modul pengembangan awal yang berfokus pada validasi algoritma pembentukan graf dinamis. Notebook ini digunakan untuk melakukan \textit{tuning} parameter ambang batas (\textit{threshold}) korelasi dan jendela waktu (\textit{window size}) sebelum diintegrasikan ke modul utama.
    
    \item \textbf{\texttt{ai\_gated\_markowitz\_complete.ipynb}} \\
    Implementasi mandiri dari strategi Baseline (AI-Gated Markowitz Klasik). Notebook ini berfungsi sebagai kontrol eksperimen untuk memastikan performa baseline terukur secara objektif tanpa bias dari modul graph.
\end{enumerate}

\subsection{Eksperimen Pembanding Lanjutan (GNN)}
Untuk menguji hipotesis efektivitas metode eksplisit \textit{feature extraction} dibandingkan dengan \textit{end-to-end deep learning}, dilakukan eksperimen tambahan menggunakan Graph Neural Networks.

\begin{enumerate}
    \item \textbf{\texttt{gnn\_enhanced\_comparison.ipynb}} \\
    Implementasi model \textit{Graph Convolutional Networks} (GCN) dan \textit{Graph Attention Networks} (GAT). Notebook ini membandingkan apakah pembelajaran representasi graf otomatis dapat mengungguli fitur graf manual. Hasil eksperimen ini digunakan dalam diskusi tesis untuk menjustifikasi pemilihan metode Random Forest + Network Features.
    
    \item \textbf{\texttt{gnn\_simplified.ipynb}} \\
    Versi ringan dari implementasi GNN yang dikembangkan untuk validasi konsep cepat tanpa ketergantungan pada library C++ yang kompleks, memastikan kode dapat dijalankan di berbagai lingkungan komputasi.
\end{enumerate}

\subsection{Eksperimen Pendahulu (Reinforcement Learning)}
Sebelum mengerucut pada pendekatan Temporal Graph, penelitian ini mengeksplorasi pendekatan \textit{Deep Reinforcement Learning} (DRL).

\begin{itemize}
    \item \textbf{\texttt{rl\_portfolio\_ultimate.ipynb}}: Implementasi agen PPO (\textit{Proximal Policy Optimization}) untuk manajemen portofolio. Kode ini dinonaktifkan dalam hasil akhir namun disimpan sebagai bukti eksplorasi metode (\textit{negative results}) yang mengarahkan penelitian pada pendekatan Graph Analytics yang lebih \textit{interpretable}.
\end{itemize}

\subsection{Ringkasan Berkas}
Daftar lengkap berkas kode sumber beserta deskripsi singkatnya disajikan dalam Tabel \ref{tab:daftar_notebook}.

\begin{table}[ht]
    \centering
    \caption{Ringkasan Artefak Kode Sumber Eksperimen}
    \label{tab:daftar_notebook}
    \renewcommand{\arraystretch}{1.3}
    \begin{tabular}{|p{5.5cm}|p{9.5cm}|}
    \hline
    \textbf{Nama Berkas Notebook} & \textbf{Deskripsi Fungsi Utama} \\ \hline
    \texttt{comparison\_baseline\_vs\newline\_temporal\_graph\_fixed.ipynb} & \textbf{(FINAL)} Analisis komparatif lengkap antara Baseline dan Temporal Graph-Enhanced Strategy beserta visualisasi hasil. \\ \hline
    \texttt{temporal\_graph\_markowitz.ipynb} & Modul pengembangan algoritma konstruksi graf dinamis dan ekstraksi fitur sentralitas. \\ \hline
    \texttt{ai\_gated\_markowitz\newline\_complete.ipynb} & Kode sumber strategi Baseline (Random Forest + Markowitz standar). \\ \hline
    \texttt{gnn\_enhanced\_comparison.ipynb} & Eksperimen pembanding menggunakan metode \textit{Graph Neural Networks} (GCN/GAT). \\ \hline
    \texttt{thesis\_experiment\_final.ipynb} & Notebook konsolidasi untuk agregasi data dan pelaporan metrik akhir. \\ \hline
    \end{tabular}
\end{table}
